% GENERATED FILE. Edit canonical lessons, then run npm run generate:books.
% canonical-source-hash: fnv1a64:f10d2c2aa87a0da6
% canonical-lessons: PA-C84-pila, PA-C84-lal, PA-C84-nila, PA-C84-kala, PA-C84-chitta

\chapter{Yellow, Red, Blue, Black, White}
\label{ch:pa-yellow-red-blue-black-white}
\hlchaptermodality{84}

\begin{chapteropening}
\textbf{By the end of this chapter:} \emph{I can say yellow, red, blue, black and white.}

Complete the last lesson of chapter 84: I can say yellow, red, blue, black and white.
\end{chapteropening}

\section[\texorpdfstring{\pa{ਪੀਲਾ} (pīlā)}{ਪੀਲਾ (pīlā)}]{\pa{ਪੀਲਾ} (pīlā) — yellow}

\label{lesson:PA-C84-pila}



\begin{quote}
Before the new one: say the Punjabi for important, then the Punjabi for green.
\end{quote}

\subsection*{You'll want to know: \pa{ਪੀਲਾ}}
\textbf{\pa{ਪੀਲਾ}} — \emph{pīlā} — \textquotedblleft{}yellow\textquotedblright{}.

\begin{cousinweb}
\textbf{\pa{ਪੀਲਾ}} (\emph{pīlā}) is a tadbhava, a word inherited from Sanskrit and reshaped by sound change. It comes from \emph{pītala}, \textquotedblleft{}yellow\textquotedblright{} (CDIAL 8233). It is built on Sanskrit \emph{pīta}, \textquotedblleft{}yellow\textquotedblright{}.
\end{cousinweb}

\begin{grammarlens}[title={What you've built}]
One more word for this chapter.
\end{grammarlens}

\subsection*{Your turn}
\begin{itemize}
\raggedright
  \item \emph{Say it:} \emph{pīlā}
  \item \emph{Say it:} \emph{pīlā}, once more
  \item \emph{Say it:} say \emph{pīlā}
  \item \emph{Recall:} say the Punjabi for important, then the Punjabi for green, then say \emph{pīlā} again
\end{itemize}

\begin{tcolorbox}[breakable,colback=teal!4,colframe=teal!35!black,title={Before you move on}]
What does \pa{ਪੀਲਾ} mean? (\textquotedblleft{}Yellow\textquotedblright{}.) Say it once more.
\end{tcolorbox}
\section[\texorpdfstring{\pa{ਲਾਲ} (lāl)}{ਲਾਲ (lāl)}]{\pa{ਲਾਲ} (lāl) — red}

\label{lesson:PA-C84-lal}



\begin{quote}
Before the new one: say the Punjabi for green, then the Punjabi for yellow.
\end{quote}

\subsection*{You'll want to know: \pa{ਲਾਲ}}
\textbf{\pa{ਲਾਲ}} — \emph{lāl} — \textquotedblleft{}red\textquotedblright{}.

\begin{cousinweb}
\textbf{\pa{ਲਾਲ}} (\emph{lāl}) is a loan, usually traced to Persian \emph{lāl}, \textquotedblleft{}red\textquotedblright{}.
\end{cousinweb}

\begin{grammarlens}[title={What you've built}]
One more word for this chapter.
\end{grammarlens}

\subsection*{Your turn}
\begin{itemize}
\raggedright
  \item \emph{Say it:} \emph{lāl}
  \item \emph{Say it:} \emph{lāl}, once more
  \item \emph{Say it:} say \emph{lāl}
  \item \emph{Recall:} say the Punjabi for green, then the Punjabi for yellow, then say \emph{lāl} again
\end{itemize}

\begin{tcolorbox}[breakable,colback=teal!4,colframe=teal!35!black,title={Before you move on}]
What does \pa{ਲਾਲ} mean? (\textquotedblleft{}Red\textquotedblright{}.) Say it once more.
\end{tcolorbox}
\section[\texorpdfstring{\pa{ਨੀਲਾ} (nīlā)}{ਨੀਲਾ (nīlā)}]{\pa{ਨੀਲਾ} (nīlā) — blue}

\label{lesson:PA-C84-nila}



\begin{quote}
Before the new one: say the Punjabi for yellow, then the Punjabi for red.
\end{quote}

\subsection*{You'll want to know: \pa{ਨੀਲਾ}}
\textbf{\pa{ਨੀਲਾ}} — \emph{nīlā} — \textquotedblleft{}blue\textquotedblright{}.

\begin{cousinweb}
\textbf{\pa{ਨੀਲਾ}} (\emph{nīlā}) is a tadbhava, a word inherited from Sanskrit and reshaped by sound change. It comes from \emph{nīla}, \textquotedblleft{}dark blue, dark green, black\textquotedblright{} (CDIAL 7563).
\end{cousinweb}

\begin{grammarlens}[title={What you've built}]
One more word for this chapter.
\end{grammarlens}

\subsection*{Your turn}
\begin{itemize}
\raggedright
  \item \emph{Say it:} \emph{nīlā}
  \item \emph{Say it:} \emph{nīlā}, once more
  \item \emph{Say it:} say \emph{nīlā}
  \item \emph{Recall:} say the Punjabi for yellow, then the Punjabi for red, then say \emph{nīlā} again
\end{itemize}

\begin{tcolorbox}[breakable,colback=teal!4,colframe=teal!35!black,title={Before you move on}]
What does \pa{ਨੀਲਾ} mean? (\textquotedblleft{}Blue\textquotedblright{}.) Say it once more.
\end{tcolorbox}
\section[\texorpdfstring{\pa{ਕਾਲਾ} (kālā)}{ਕਾਲਾ (kālā)}]{\pa{ਕਾਲਾ} (kālā) — black}

\label{lesson:PA-C84-kala}



\begin{quote}
Before the new one: say the Punjabi for red, then the Punjabi for blue.
\end{quote}

\subsection*{You'll want to know: \pa{ਕਾਲਾ}}
\textbf{\pa{ਕਾਲਾ}} — \emph{kālā} — \textquotedblleft{}black\textquotedblright{}.

\begin{cousinweb}
\textbf{\pa{ਕਾਲਾ}} (\emph{kālā}) is a tadbhava, a word inherited from Sanskrit and reshaped by sound change. It comes from \emph{kāla}, \textquotedblleft{}black, dark-blue\textquotedblright{} (CDIAL 3083). CDIAL derives the Sanskrit word itself, through \emph{*kāḍa}, from Dravidian.
\end{cousinweb}

\begin{grammarlens}[title={What you've built}]
One more word for this chapter.
\end{grammarlens}

\subsection*{Your turn}
\begin{itemize}
\raggedright
  \item \emph{Say it:} \emph{kālā}
  \item \emph{Say it:} \emph{kālā}, once more
  \item \emph{Say it:} say \emph{kālā}
  \item \emph{Recall:} say the Punjabi for red, then the Punjabi for blue, then say \emph{kālā} again
\end{itemize}

\begin{tcolorbox}[breakable,colback=teal!4,colframe=teal!35!black,title={Before you move on}]
What does \pa{ਕਾਲਾ} mean? (\textquotedblleft{}Black\textquotedblright{}.) Say it once more.
\end{tcolorbox}
\section[\texorpdfstring{\pa{ਚਿੱਟਾ} (chiṭṭā)}{ਚਿੱਟਾ (chiṭṭā)}]{\pa{ਚਿੱਟਾ} (chiṭṭā) — white}

\label{lesson:PA-C84-chitta}



\begin{quote}
Before the new one: say the Punjabi for blue, then the Punjabi for black.
\end{quote}

\subsection*{You'll want to know: \pa{ਚਿੱਟਾ}}
\textbf{\pa{ਚਿੱਟਾ}} — \emph{chiṭṭā} — \textquotedblleft{}white\textquotedblright{}.

\begin{cousinweb}
\textbf{\pa{ਚਿੱਟਾ}} (\emph{chiṭṭā}) is a tadbhava, a word inherited from Sanskrit and reshaped by sound change. It comes from \emph{śvitra}, \textquotedblleft{}white\textquotedblright{} (CDIAL 12772).
\end{cousinweb}

\begin{grammarlens}[title={What you've built}]
One more word for this chapter.
\end{grammarlens}

\subsection*{Your turn}
\begin{itemize}
\raggedright
  \item \emph{Say it:} \emph{chiṭṭā}
  \item \emph{Say it:} \emph{chiṭṭā}, once more
  \item \emph{Say it:} say \emph{chiṭṭā}
  \item \emph{Recall:} say the Punjabi for blue, then the Punjabi for black, then say \emph{chiṭṭā} again
\end{itemize}

\begin{tcolorbox}[breakable,colback=teal!4,colframe=teal!35!black,title={Before you move on}]
What does \pa{ਚਿੱਟਾ} mean? (\textquotedblleft{}White\textquotedblright{}.) Say it once more.
\end{tcolorbox}
